% id: 88-08
% book: 베이직쎈 확률과 통계 · 05 확률변수와 확률분포
% section: 유형 049 확률질량함수
% figure: none
% answer: 
% answer_source: 
% variant_level: 0 (원본 전사)
% 이 파일은 build.mjs 가 items.json 에서 만든다. 고칠 때는 items.json 을 고친다.
\smash{\raisebox{1.5mm}{\dmkichul{베이직쎈 확률과 통계 88쪽 08번}}}
$1$부터 $10$까지의 자연수가 각각 하나씩 적힌 10개의 공이 들어 있는 상자에서 임의로 2개의 공을 동시에 꺼낼 때, 꺼낸 소수가 적힌 공의 개수를 확률변수 $X$라 하자. $X$의 확률질량함수가\par\smallskip\dispeq{$\displaystyle \mathrm{P}(X=x)=\dfrac{{}_p\mathrm{C}_x\cdot{}_q\mathrm{C}_{2-x}}{{}_{10}\mathrm{C}_r}\ (x=0{,}\mkern3mu\mathopen{}1{,}\mkern3mu\mathopen{}2)$}\par\smallskip\noindent 일 때, 상수 $p$, $q$, $r$의 값은?
\begin{choicesii}{$p=4$, $q=6$, $r=2$}{$p=4$, $q=6$, $r=4$}{$p=5$, $q=5$, $r=2$}{$p=6$, $q=4$, $r=4$}{$p=6$, $q=4$, $r=2$}\end{choicesii}
